% Figure: a simply-supported beam under a central point load, with its
% shear diagram (step) and moment diagram (triangle) stacked below.
\begin{figure}[htbp]
\centering
\begin{tikzpicture}[>=latex, x=1.0cm, y=1.0cm]
  \def\L{6.0}
  % ==== beam ====
  \begin{scope}
    \draw[engblue, very thick] (0,0) -- (\L,0);
    % supports
    \draw[thick] (0,0) -- (-0.3,-0.45) -- (0.3,-0.45) -- cycle;
    \draw[thick] (\L,-0.18) circle (0.12);
    \draw[thick] (\L-0.3,-0.3) -- (\L+0.3,-0.3);
    % central point load
    \draw[->, very thick, engred] (\L/2,1.2) -- (\L/2,0.1);
    \node[engred, anchor=south, font=\scriptsize] at (\L/2,1.2) {$P$};
    % reactions
    \node[font=\scriptsize, anchor=north east] at (0,-0.45) {$P/2$};
    \node[font=\scriptsize, anchor=north west] at (\L,-0.45) {$P/2$};
  \end{scope}
  % ==== shear diagram ====
  \begin{scope}[yshift=-2.6cm]
    \draw[->, thick] (-0.3,0) -- (\L+0.4,0) node[right, font=\scriptsize] {$x$};
    \node[font=\scriptsize, anchor=south east] at (0,0.9) {$V(x)$};
    % +P/2 from 0 to L/2, then -P/2
    \draw[engblue, very thick] (0,0.8) -- (\L/2,0.8) -- (\L/2,-0.8) -- (\L,-0.8);
    \draw[engblue, very thick, dotted] (0,0) -- (0,0.8);
    \draw[engblue, very thick, dotted] (\L,-0.8) -- (\L,0);
    \node[font=\scriptsize, anchor=south] at (\L/4,0.8) {$+P/2$};
    \node[font=\scriptsize, anchor=north] at (3*\L/4,-0.8) {$-P/2$};
  \end{scope}
  % ==== moment diagram ====
  \begin{scope}[yshift=-5.2cm]
    \draw[->, thick] (-0.3,0) -- (\L+0.4,0) node[right, font=\scriptsize] {$x$};
    \node[font=\scriptsize, anchor=south east] at (0,1.1) {$M(x)$};
    % triangle peaking at L/2 with PL/4
    \draw[engred, very thick] (0,0) -- (\L/2,1.0) -- (\L,0);
    \draw[enggray, dashed, thin] (\L/2,0) -- (\L/2,1.0);
    \node[engred, font=\scriptsize, anchor=south] at (\L/2,1.0) {$M_{\max}=PL/4$};
  \end{scope}
\end{tikzpicture}
\caption[Shear and moment diagrams: central point load]{The shear and
moment diagrams of a simply-supported beam under a central point load,
the case worth carrying in memory. The reactions are $P/2$ at each
support by symmetry. The shear is constant at $+P/2$ from the left
support to the load, drops by $P$ at the load (a concentrated load makes
a jump in shear equal to its size), and continues at $-P/2$ to the right
support. The moment rises linearly with slope equal to the shear,
$dM/dx=V$, reaching its peak $PL/4$ exactly where the shear crosses zero
under the load, then falls back to zero at the right support. The kink in
the moment diagram under the load, with no jump, is the signature of a
concentrated transverse load; a concentrated applied couple would instead
put a jump in the moment diagram and leave the shear unchanged. Reading
these two rules off the load lets the engineer sketch both diagrams
without integrating.}
\label{fig:vol03ch02:shear-moment}
\end{figure}
